%% ma: L11-B27-0008
%% lop: 11
%% chuong: 7
%% bai: 27
%% dang: 11.27.6
%% loai: TLN
%% muc: VD
%% dap_an: "24"
%% nguon: "(NBV)-ĐỀ SỐ 5-MỖI NGÀY 1 ĐỀ THI 2025.pdf (D05, MỖI NGÀY 1 ĐỀ - Đề số 5), Phần III Câu 4"
%% kien_thuc: "Bài 27 (thể tích khối lăng trụ, khối chóp; định lí Pythagore tính chiều cao chóp)"
%% trang_thai: chinh-thuc
%% ghi_chu: "Giáo viên đã duyệt 10/10/2026. Đề gốc không nói rõ 3 cm là gì; theo lời giải gốc là độ dài cạnh bên của khối chóp, đã thêm chú thích trong đề. Hình dựng lại không theo tỉ lệ chính xác"
\begin{ex}
Người ta thiết kế một thiết bị kim loại có dạng như hình vẽ (giá tiền mua kim loại là $2500$ đồng$/\mathrm{cm}^3$). Thiết bị gồm $2$ phần, phần dưới là khối lăng trụ tứ giác đều, phần trên là khối chóp tứ giác đều ($3$ cm là độ dài cạnh bên của khối chóp). Số tiền mua kim loại dùng để làm thiết bị đó là bao nhiêu nghìn đồng (làm tròn kết quả đến hàng đơn vị)?
\begin{center}
\begin{tikzpicture}[scale=.8]
  \path (0,0) coordinate (A) (3,0) coordinate (B) (4,1) coordinate (C) (1,1) coordinate (D);
  \foreach \P in {A,B,C,D} \path (\P)+(0,2.25) coordinate (\P');
  \coordinate (S) at (2,6.72);
  \draw (A)--(B)--(C) (A)--(A') (B)--(B') (C)--(C') (A')--(B')--(C') (S)--(A') (S)--(B') (S)--(C');
  \draw[khuat] (A)--(D)--(C) (D)--(D') (A')--(D')--(C') (S)--(D');
  \draw[<->] (-.6,0)--(-.6,2.25) node[midway,left]{$1{,}5$ cm};
  \path (A) -- (B) node[midway,below]{$2$ cm};
  \path (B) -- (C) node[midway,below right]{$2$ cm};
  \path (S) -- (C') node[midway,right]{$3$ cm};
\end{tikzpicture}
\end{center}
\shortans{24}
\loigiai{Lăng trụ tứ giác đều đáy cạnh $2$ cm, cao $1{,}5$ cm: $V_1=2^2\cdot1{,}5=6$ ($\mathrm{cm}^3$).
Chóp tứ giác đều đáy cạnh $2$ cm, cạnh bên $3$ cm: nửa đường chéo đáy bằng $\sqrt2$ cm nên chiều cao $h=\sqrt{3^2-(\sqrt2)^2}=\sqrt7$ cm và $V_2=\dfrac13\cdot2^2\cdot\sqrt7=\dfrac{4\sqrt7}{3}$ ($\mathrm{cm}^3$).
Số tiền: $2500\left(6+\dfrac{4\sqrt7}{3}\right)\approx23\,819$ đồng $\approx24$ nghìn đồng.}
\end{ex}
